\clearpage
\phantomsection
\section*{Yang--Mills: state of the problem}
\addcontentsline{toc}{section}{Yang--Mills: state of the problem}
\label{ym:sec:status}
\textit{Construction assessment as of 23 November 2026; finite-torus results updated 25 November 2026 (UTC).}

\paragraph{The mathematical target.}
For every compact simple gauge group, the Jaffe--Witten problem asks for a
nontrivial quantum Yang--Mills theory on $\mathbb R^4$ with the required
axiomatic and short-distance properties, and a positive spectral gap above
the vacuum \cite{JaffeWitten2006}. Constructing a theory and proving a gap
are parts of one problem; there is no separate prize for the existence part.

\paragraph{The assessed construction.}
The Yang--Mills group's assessment of 10 November supports the October
construction as a scaling limit of four-dimensional Wilson lattice gauge
theory, with the properties stated in Theorem~\ref{ym:thm:existence}
\cite{OpenAIYM2026,YMGroup2026}. It establishes neither a positive mass gap
nor independence of the extracted infinite-volume limit. Uniqueness of the
vacuum within a selected reconstructed phase does not establish uniqueness
of the theory obtained by extraction. The October checks also do not
establish agreement with perturbation theory to all orders
\cite{YMCheckScope2026}. The group's current listing marks the submissions
of 4 November and 15 November as work in progress; listing a submission is
not an assessment of its claims \cite{YMGroupCurrent2026}.

\paragraph{The conditional route developed here.}
Section~\ref{ym:sec:lattice} develops the finite-torus defect and center-twist
arguments. On every finite torus the transfer operator is nonnegative and
trace class (Proposition~\ref{ym:sc-prop:positivity}); the purity, defect and
conditional-gap propositions are stated in ratio form, so that they include
the rank-one transfer operator at $\beta=0$; and the defect recurrence is
a theorem (Theorem~\ref{ym:thm:recurrence}, with $c_1=3740$ and
$\eta_*=\frac1{7480}$). The finite-volume steps and the recurrence have been reviewed
independently; the strong-coupling seed (Proposition~\ref{ym:prop:strong-seed})
was proposed in that review. The general formal argument takes spectral inputs. For the
four-dimensional $\mathbb Z_2$ Wilson model the trace formula, positivity,
seed and finite-volume conclusion have now been proved together and
independently checked; see Section~\ref{ym:sec:finite-guide}. They are not part of the assessed formal construction.

At strong coupling, $\beta\le\frac{1}{7480\cdot192\,n^4}$, the chain closes
and gives a lattice gap of at least $\frac{\log2}{n}$ on every dyadic torus;
that regime already has a gap by convergent expansions. The open estimate in
this chain is Conjecture~\ref{ym:conj:defect}: a small defect on a box of
bounded physical size, with the same bare parameters on both boxes. Even
assuming it, the construction still requires compatible couplings across
volumes, estimates uniform in the cutoff for renormalized correlations, and
control of the infinite-volume limit; the last is stated as an explicit
hypothesis after Eq.~\eqref{ym:finite-correlation-bound}. A finite-volume gap
or an exact check on a finite lattice does not establish
Eq.~\eqref{ym:gap-statement} for the continuum theory. No positive continuum
gap is claimed here.

\paragraph{Lattice measurements.}
These are finite numerical checks, not estimates in the proof.
\begin{itemize}
\item \emph{Defect.} For $SU(2)$, $\Delta_\beta(2)$ is measured for $\beta\le2.4$,
 and $\Delta_\beta(4)$ as increments from $\beta=1$. A truncated
 strong-coupling formula matches the measurements for $\beta\le0.8$
 (Section~\ref{ym:num-defect}); it is not a controlled expansion. No absolute
 defect has been measured at weak coupling.
\item \emph{A necessary condition for a seed.} A seed at side $n$ requires
 $n\,\mathsf m(n)\ge10.9994$ (Eq.~\eqref{ym:num-necessary}): the box must be
 at least about eleven of its own correlation lengths across.
\item \emph{$U(1)$.} At $\beta=1.6$ the defect at $n=2$ is $88.48(3)$;
 at $n=4$ the measured \emph{increment} from $\beta=0.6$ is $80.3(1)$.
 Neither is evidence for a small seed.
\item \emph{Glueball.} At $\beta=2.5$ the $0^{++}$ mass is $0.730(11)$ on
 $20^4$ and $0.691(17)$ on $16^4$, against $0.723(23)$ and $0.660(22)$ in
 \cite{Teper1998}. At $L=8,10,12$ neither designated fit window passes; mixing with
 torelon pairs is a hypothesis there. Single-torelon energies give
 $a\sqrt\sigma=0.183$--$0.189$, against $0.1834(26)$ and $0.1881(28)$
 (Section~\ref{ym:num-glueball}).
\item \emph{Topology.} The count of vortex sheets tracks confinement on
 $8^4$--$16^4$ for one gauge-fixing prescription
 (Section~\ref{ym:num-topology}). This is an association, not an order
 parameter.
\end{itemize}

\paragraph{The prize status.}
Clay continues to list the problem as unsolved \cite{ClayYMStatus2026}.
Its rules require a complete solution, publication in a qualifying outlet,
at least two years after publication, and general acceptance before
consideration \cite{ClayPrizeRules2018}. The assessment of the October
construction is not a Clay award decision. A future infrared result would
need to be checked against the complete problem, with appropriate credit
for the construction and any other indispensable prior work. There is no
mathematically meaningful percentage of the prize established by the
number of calculations or lemmas completed.

\paragraph{Reading route.}
For a route from holonomy through twist sectors to a finite-size spectral
bound, start with Section~\ref{ym:sec:finite-guide}
(p.~\pageref{ym:sec:finite-guide}). For the QFT background begin with
relativistic fields in Section~\ref{qft:sec:scalar}
(p.~\pageref{qft:sec:scalar}) and quantization in
Section~\ref{qft:sec:quantization}
(p.~\pageref{qft:sec:quantization}), using Srednicki
\cite{MSrednicki2007} for the QFT background. Continue with the lattice
definitions and scaling limit in Section~\ref{ym:sec:lattice}
(p.~\pageref{ym:sec:lattice}); then separate the existence theorem from
the spectral question and the explicitly conditional infrared programme.
Section~\ref{ym:sec:source-map}
(p.~\pageref{ym:sec:source-map}) locates the statements in the formal source.
\clearpage
